\section{Data Resources, Application Space, and Evaluation Logic}

Recent remote sensing review articles in RSE, TGRS, and ISPRS tend to avoid treating methods, datasets, and applications as isolated lists. Instead, they first describe the sensing and data regime, then organize methodological families around the downstream questions that remote sensing users actually ask, and finally return to limitations, scale effects, and evaluation practice \citep{ma2019dl_meta_isprs,yuan2020env_dl_rse,weiss2020agri_meta_rse,paoletti2019hsi_review,li2020object_detection_isprs,survey1,xiao2025foundation_survey}. This organization is especially appropriate for satellite embeddings because an embedding is not a complete solution by itself. It is an intermediate representation whose usefulness depends on the sensor stream from which it was learned, the spatial and temporal scale at which it is exposed, and the type of downstream evidence that a user wants to recover. A pooled vector can be excellent for tile classification but insufficient for shoreline delineation; a dense feature grid can support segmentation but may be expensive or unavailable for some products; an annual precomputed layer can be ideal for wide-area screening but poorly matched to event-specific monitoring. The application section therefore has to do more than enumerate tasks. It must explain how the embedding form constrains the mapping problem.

The first organizing dimension is the data regime. Most current satellite embeddings are produced from one of four recurring sources: medium-resolution optical archives such as Sentinel-2 and Landsat, optical-SAR combinations centered on Sentinel-1 and Sentinel-2, broader multimodal stacks that add elevation, weather, labels, text, or metadata, and precomputed embedding products that have already collapsed a large archive into a queryable feature layer. These regimes create different forms of transfer. Optical-only models usually retain strong spectral and spatial cues for land-cover mapping, but they may fail under cloud, season, or sensor changes if temporal handling is weak. Optical-SAR and multimodal models try to absorb those changes during representation learning, but their practical value depends on whether all required inputs are available in the target region and period. Precomputed products shift the problem again: they reduce inference cost and make wide-area sampling simple, yet they also fix the upstream temporal summary, native scale, and production assumptions. This is why the same model family can behave very differently when evaluated at 10\,m, 30\,m, 500\,m, 1\,km, or 0.25$^\circ$ resolution. In an embedding-centered review, native scale is not a footnote; it is part of the representation.

The second organizing dimension is the downstream target. Classification and mapping remain the most common entry point because they are compatible with almost every embedding family. A single pooled vector per tile is often enough to support land-cover labels, crop categories, forest condition classes, urban typologies, or simple ecological strata. This makes classification a useful baseline, but it is also an incomplete test: a model can separate dominant classes while losing continuous mixture information or local boundary structure. Regression tasks address that gap by asking whether the embedding preserves smoothly varying quantities such as built-up fraction, vegetation fraction, biomass proxies, canopy height, crop-yield indicators, or environmental indices. Dense prediction tasks push further by asking whether the representation still contains spatially localized information after tokenization, pooling, product resampling, or temporal summarization. Water extraction, land-cover segmentation, flood mapping, and mask reconstruction are therefore not merely additional applications; they probe whether the embedding is a spatial representation rather than only a semantic descriptor. Finally, vision-language and cross-modal embeddings introduce retrieval and open-vocabulary search as a distinct mode of use. These tasks are promising, but they should be interpreted separately from supervised mapping because they measure semantic alignment and promptability rather than only geographic fidelity.

The third organizing dimension is operational fit. Remote sensing users rarely choose an embedding only because it has the highest number in a table. They also need to know whether the representation can be generated for the desired sensor, season, and region; whether it exposes pooled or dense outputs; whether it can be sampled as a product or requires GPU inference; whether runtime and memory scale with area; and whether the temporal meaning of the output is annual, seasonal, monthly, or scene-specific. Recent foundation-model surveys have emphasized the same concern from a broader model perspective, but embedding selection makes it even sharper: two systems can both be called foundation models while exposing very different reusable objects \citep{survey1,survey3,survey4,xiao2025foundation_survey}. The review therefore treats applications as representation tests. Classification tests semantic separability, regression tests continuous structure, binary extraction tests operational thresholdability, segmentation tests local spatial separability, and retrieval tests cross-modal alignment. This evaluation logic leads directly to the compact benchmark in the next section.
